\section{Context Coherence Score and HCPC-v$2$}
\label{sec:method}

This section formalises the diagnostic and the deployment policy.
Both are deliberately small.

\subsection{Context Coherence Score (CCS)}
\label{sec:method:ccs}

Let $\mathcal{C} = \{c_1, \ldots, c_k\}$ denote the top-$k$ retrieved
chunks for a query, and let $\mathbf{e}_i \in \mathbb{R}^d$ denote
their L$2$-normalised embeddings. The Context Coherence Score is:

\begin{equation}
\textrm{CCS}(\mathcal{C}) \;=\;
  \mu_{i<j}\!\bigl[\mathbf{e}_i \cdot \mathbf{e}_j\bigr]
  \;-\;
  \sigma_{i<j}\!\bigl[\mathbf{e}_i \cdot \mathbf{e}_j\bigr],
\end{equation}

where the mean and standard deviation are taken over the $\binom{k}{2}$
off-diagonal pairwise cosine similarities. The subtraction term is
load-bearing: a set with high mean similarity but high variance --- two
tight sub-clusters joined by no bridge --- has the same mean as a set
that is uniformly mid-similarity, and the latter is what the generator
needs. CCS rewards the latter.

CCS is well-defined for any $k \geq 2$, requires only the embeddings
already produced by the retrieval step, runs in $O(k^2 d)$ time,
and is generator-free. For $k = 3$ (our default) it is $3 d$
floating-point operations per query --- negligible relative to the
generator forward pass.

\paragraph{Decision policy.}
We treat CCS as a deployment-time gate. Given a similarity threshold
$\tau \in [-1, 1]$, the bare CCS gate is:

\begin{quote}\small\ttfamily
docs, sims = retriever.retrieve(query)\\
ccs = compute\_ccs(docs, embedder)\\
if ccs $\geq$ tau:\\
\hspace*{2em}return docs                       \# coherent — pass to generator\\
else:\\
\hspace*{2em}return refine(docs)               \# incoherent — apply HCPC-v1
\end{quote}

The bare CCS gate already recovers most of HCPC-v$2$'s gain
(\S\ref{sec:robustness}, Table~\ref{tab:robustness_summary}). We deploy it as the
\emph{decision policy}, and use HCPC-v$2$ as the implementation that
adds two additional safeguards described next.

\subsection{HCPC-v$2$: protected ranks $+$ sub-chunk merge-back}
\label{sec:method:hcpc_v2}

HCPC-v$2$ extends the bare gate with two additional rules tuned on
small held-out sets:

\paragraph{Rank-protected top-$k$.}
Even when the gate decides to refine, the top-$2$ retrieved chunks
(by raw query similarity) are \emph{never} touched. These are the
chunks most likely to carry the answer span; refining them risks the
exact failure mode the gate is meant to prevent.

\paragraph{Sub-chunk merge-back.}
Chunks that pass the gate-and-rank filter are subdivided into
$256$-token sub-chunks, each rescored by the cross-encoder, and only
the top sub-chunk per parent is retained. The result is that
refinement, when it happens, replaces a passage with its most
discriminating fragment rather than a fragment from elsewhere in the
corpus. We sweep the sub-chunk size in $\{128, 256, 512\}$ and report
sensitivity in \S\ref{sec:robustness}.

\paragraph{Hyperparameters.}
The coherence threshold $\tau = 0.50$ and the similarity / cross-encoder
thresholds for the dual-AND gate ($\sigma_{\textrm{sim}} = 0.45$,
$\sigma_{\textrm{ce}} = -0.20$) are tuned on a single held-out
set of $50$ SQuAD queries and held fixed across every experiment in the
paper. We sweep $\tau \in \{0.4, 0.5, 0.6\}$ in
Appendix~\ref{app:hyperparams} and confirm the policy is stable
across a $\pm 0.1$ window around the operating point.

\paragraph{Cost.}
Per query, HCPC-v$2$ adds: one CCS computation
($O(k^2 d) = $ a few thousand FLOPs); when the gate fires, $O(k)$
cross-encoder calls ($\sim 5$ ms each on CPU); the optional
sub-chunk pass adds $O(k \cdot \textrm{n\_subs})$ further cross-encoder
calls. Total marginal latency over the baseline pipeline: typically
$<50$~ms, dominated by the cross-encoder. For comparison, Self-RAG's
reflection-token mechanism roughly doubles end-to-end latency.

\subsection{Why we ship the bare gate as the deployment policy}
\label{sec:method:why_bare}

A reasonable reviewer concern: HCPC-v$2$ has three moving parts
(coherence gate, rank protection, sub-chunk merge-back), and the
ablation should isolate which part is doing the work. We construct
the bare CCS gate (CCSGate, no rank protection, no sub-chunks) and
report its performance in the top-$k$ ablation
(Table~\ref{tab:robustness_summary}, ``gate~rec.'' column). On SQuAD the bare gate
recovers $0.067$ of HCPC-v$2$'s $0.062$ recovery at $k=2$ --- nearly
identical. The protected-rank rule contributes the bulk of the
remaining gain only at larger $k$ (where there are more low-rank
chunks at risk), and the sub-chunk merge-back is a refinement-quality
improvement rather than a safety mechanism.

The implication: \textbf{the gate decision is the active ingredient.}
We ship it as the recommended deployment policy under the name
\texttt{CCSGate}; HCPC-v$2$ is the production implementation with the
two extra safeguards for $k > 3$ deployments. Both are released as a
$\texttt{pip install context-coherence}$ package and as a drop-in
LangChain \texttt{CoherenceGatedRetriever}.
